% META: {"id": "1NSI-TC-EX-002", "chapitre": "1NSI-TYPES-CONSTRUITS", "type_objet": "exercice", "capacites": ["1NSI-TYPES-CONSTRUITS-C1"], "parcours": 2, "competences": ["analyser"], "duree_min": 8, "mode_creation": "transposition", "sources_inspiration": [], "status": "approved", "capacites_codes": ["C1"], "methodes": ["M1"], "fichier_tex": "chapitres/1NSI-TYPES-CONSTRUITS/exercices/1NSI-TC-EX-002.tex", "corrige_tex": "chapitres/1NSI-TYPES-CONSTRUITS/corriges/1NSI-TC-CO-002.tex"}
% BEGIN-TRACE
% station = ("Tunis", 36.8, 10.2)
% print(type(station))
% print(station[0])
% print(len(station))
% EXPECTED
% <class 'tuple'>
% Tunis
% 3
% END-TRACE
\begin{exercice}{1NSI-TC-EX-002}{1}{5}
On donne le tuple suivant :
\begin{python}
station = ("Tunis", 36.8, 10.2)
\end{python}
\begin{enumerate}
  \item Donner le type de \lstinline{station} et le nombre de ses éléments.
  \item Donner la valeur de \lstinline{station[0]} et de \lstinline{station[2]}.
  \item Peut-on écrire \lstinline{station[1] = 37.0} ? Justifier.
\end{enumerate}
\end{exercice}
